Fix \(P_{\mathrm{obs}}\). If the coordinate cell is negative for \((a,y)\), \(\text{def:terminal-qe-test}\) and \(\text{def:symbolic-observed-partition}\) make the formula
\[
\exists\theta^0\,\exists\theta^1
\left[
\mathsf C_{\mathcal G}(\theta^0;P_{\mathrm{obs}})
\wedge\mathsf C_{\mathcal G}(\theta^1;P_{\mathrm{obs}})
\wedge q_{a,y}(\theta^0)\neq q_{a,y}(\theta^1)
\right]
\tag{1}
\]
true over \(\mathbb R\). Choose a satisfying pair. The converse direction of \(\text{thm:response-fiber-equivalence}\) decodes it to \(\mathcal M_0,\mathcal M_1\in\mathfrak F(\mathcal G,P_{\mathrm{obs}})\) with unequal \(Q_a(y)\) values.

If instead the direct target cell is negative at \(\mathbf s_Y=(s_Y(y))_y\), the formula
\[
\exists\theta^0\,\exists\theta^1
\left[
\mathsf C_{\mathcal G}(\theta^0;P_{\mathrm{obs}})
\wedge\mathsf C_{\mathcal G}(\theta^1;P_{\mathrm{obs}})
\wedge t_{\mathbf s_Y}(\theta^0)\neq t_{\mathbf s_Y}(\theta^1)
\right]
\tag{2}
\]
is true. Apply the converse direction of \(\text{thm:response-fiber-equivalence}\) to both response vectors. It gives
\[
\mathcal M_0,\mathcal M_1\in\mathfrak F(\mathcal G,P_{\mathrm{obs}}),
\qquad
P_{\mathcal M_0}(\mathcal O)=P_{\mathcal M_1}(\mathcal O)=P_{\mathrm{obs}}.
\tag{3}
\]
By the definition of \(t_{\mathbf s_Y}\), coordinate preservation in response-fiber equivalence, finite summation, and \(\text{def:causal-query}\), (2) becomes
\[
\tau_{s_Y}^{\mathcal M_0}=t_{\mathbf s_Y}(\theta^0)
\neq t_{\mathbf s_Y}(\theta^1)=\tau_{s_Y}^{\mathcal M_1}.
\tag{4}
\]
The inclusion \(\mathfrak M_{+}^{\mathrm{rm}}(\mathcal G)\subseteq\mathfrak M_{\mathrm{lm}}(\mathcal G)\) follows from \(\text{prop:positive-subclass}\). Thus either pair also lies in the corresponding ambient observed-law set. Constancy of the relevant query on that larger set would contradict (1) or (4), proving ambient nonidentification. The choices in (1)--(2) use ordinary real semantics, not an effective encoding. The exact interfaces are required only to output represented points.